% from ega4-2-fr.tex
\oldpage[IV-2]{120}

\item[(ii)] L'ensemble $U$ des $x\in X$ tels que $\mathcal F$ vérifie $(S_2)$
au point $x$ (5.7.2) est ouvert dans $X$ et on a
$\operatorname{codim}(X-U,X)\geq2$.
\end{enumerate}
\end{corollary}

(i) On sait (5.9.11) que $\mathcal F'$ est $Z^{(2)}$-clos, et en outre
$\operatorname{Supp}(\mathcal F')=X$, car les points maximaux de $X$
appartiennent à $X-Z^{(2)}$ et en ces points
$\mathcal F'_x=\mathcal F_x\neq0$, donc le support de $\mathcal F'$ est dense
dans $X$, et comme $\mathcal F'$ est cohérent,
$\operatorname{Supp}(\mathcal F')$ est fermé, donc égal à $X$. Il reste à voir
que $\mathcal F'$ n'a pas de cycles premiers associés de codimension 1. Mais
si $x\in\operatorname{Ass}(\mathcal F')$ et
$\dim(\mathcal F'_x)=\dim(\mathcal O_x)=1$, on a $x\in X-Z^{(2)}$, donc,
puisque $\mathcal F'_x=\mathcal F_x$, on aurait
$x\in\operatorname{Ass}(\mathcal F)$, contrairement à l'hypothèse, ce qui
achève de prouver (i).

(ii) On a $Z^{(n)}(X_x)=Z^{(n)}(X)\cap X_x$, avec les notations de (5.9.6),
compte tenu de (I, 2.4.2); d'autre part, l'hypothèse que $\mathcal F$ n'a pas
de cycles premiers associés de codimension 1 entraîne la même hypothèse pour
$\widetilde{\mathcal F}_x$; pour que $\widetilde{\mathcal F}_x$ vérifie
$(S_2)$, il faut et il suffit donc, en vertu de (5.10.14), que
$\widetilde{\mathcal F}_x$ soit $Z^{(2)}(X_x)$-clos; l'assertion (ii) résulte
donc de (5.10.6) et de (5.10.10, (iii)).

\begin{proposition}[5.10.16]
\label{IV.2.s5.10.16-fr}
Soient $X$ un préschéma localement noethérien,
\[
X'=\operatorname{Spec}(\mathcal H^0_{X/Z^{(2)}}(\mathcal O_X))
\]
sa $Z^{(2)}$-clôture, $g:X'\to X$ le morphisme structural. Supposons que $X$
n'ait pas de cycle premier associé de codimension 1.
\begin{enumerate}
\item[(i)] Pour qu'en un point $x\in X$, $X$ vérifie $(S_2)$, il faut et il
suffit que le morphisme $X'_x\to X_x$ déduit de $g$ (notations de (5.10.12))
soit un isomorphisme. Cette condition est toujours vérifiée si
$\operatorname{codim}(\overline{\{x\}},X)\leq1$.
\item[(ii)] Supposons de plus que $g$ soit un morphisme fini (voir dans
(5.11.2) des conditions suffisantes pour qu'il en soit ainsi). Alors
l'ensemble $U$ des points où $X$ vérifie $(S_2)$ est ouvert et
$\operatorname{codim}(X-U,X)\geq2$; en outre $U$ est le plus grand ouvert de
$X$ tel que la restriction $g^{-1}(U)\to U$ de $g$ soit un isomorphisme.
\item[(iii)] Sous les mêmes hypothèses que dans (ii), $X'$ satisfait à $(S_2)$
et pour tout $x'\in X'$ tel que\linebreak
$\operatorname{codim}(\overline{\{x'\}},X')\leq1$, le point $x=g(x')$ est tel
que
\[
\operatorname{codim}(\overline{\{x\}},X)
=\operatorname{codim}(\overline{\{x'\}},X').
\]
\item[(iv)] Les hypothèses étant celles de (ii), soit $\mathcal F$ un
$\mathcal O_X$-Module cohérent de support $X$, tel que
$\mathcal H^0_{X/Z^{(2)}}(\mathcal F)$ soit cohérent; alors le
$\mathcal O_{X'}$-Module $\mathcal F'$ tel que
$g_*(\mathcal F')=\mathcal H^0_{X/Z^{(2)}}(\mathcal F)$ est cohérent et
vérifie $(S_2)$, et son support est une réunion de composantes irréductibles de
$X'$.
\end{enumerate}
\end{proposition}

Les assertions (i) et (ii) sont insérées pour mémoire, ayant déjà été
démontrées en substance: (i) résulte en effet de (5.10.12, (i)) et de
(5.10.14), et (ii) est un cas particulier de (5.10.15, (ii)).

Démontrons (iii); posons $x=g(x')$; comme $g$ est fini, il en est de même du
morphisme $X'_x\to X_x$, donc
\[
\dim(\mathcal O_{X',x'})\leq\dim(X'_x)\leq\dim(X_x)
=\dim(\mathcal O_{X,x})
\]
en vertu de (5.4.1). Supposons d'abord que
$\dim(\mathcal O_{X',x'})\leq1$ et montrons qu'alors
$\dim(\mathcal O_{X,x})\leq1$. Sinon, on aurait $x\in Z^{(2)}$, donc en vertu
de (5.10.12, (ii)) appliqué à $\mathcal O_X$ on aurait
$\operatorname{prof}(\mathcal O_{X',x'})\geq2$, ce qui est absurde. On a donc
$x\in X-Z^{(2)}$, et par suite $\mathcal O_{X',x'}$ est isomorphe à
$\mathcal O_{X,x}$ (5.10.12, (i)), d'où
$\dim(\mathcal O_{X',x'})=\dim(\mathcal O_{X,x})$. En outre, comme $X$ n'a
pas de cycles

